给$B$求导时，值域用法联络，两个输入槽用切联络。具体定义
\[
 (\nabla B)(X;Y,Z)=\nabla_X^\perp(B(Y,Z))
 -B(\nabla_XY,Z)-B(Y,\nabla_XZ).
\]
这是$T^*M\otimes\operatorname{Sym}^2T^*M\otimes NM$截面；若只对$B$的分量作普通偏导，会丢失三组联络项。法曲率取
$R^\perp(X,Y)\xi=\nabla_X^\perp\nabla_Y^\perp\xi-\nabla_Y^\perp\nabla_X^\perp\xi-\nabla_{[X,Y]}^\perp\xi$。

\Needspace{20\baselineskip}
\begin{theorem}[三组相容方程]\label{thm:gauss-codazzi-ricci}
对切向量场$X,Y,Z,W$和法向场$\xi,\eta$，有
\begin{align*}
 g(R(X,Y)Z,W)&=\bar g(\bar R(X,Y)Z,W)\\
 &\quad+\bar g(B(X,W),B(Y,Z))-\bar g(B(Y,W),B(X,Z)),\\
 (\bar R(X,Y)Z)^\perp&=(\nabla B)(X;Y,Z)-(\nabla B)(Y;X,Z),\\
 \bar g(R^\perp(X,Y)\xi,\eta)&=\bar g(\bar R(X,Y)\xi,\eta)
       +g([A_\xi,A_\eta]X,Y).
\end{align*}
其中$[A_\xi,A_\eta]=A_\xi A_\eta-A_\eta A_\xi$。三式依次为Gauss、Codazzi、Ricci方程；末式不是Ricci张量的缩并公式。
\end{theorem}
\begin{proof}
在$\bar\nabla_X\bar\nabla_YZ$中先代入Gauss公式，再分别对切、法两部分求导，得到
\begin{align*}
 (\bar\nabla_X\bar\nabla_YZ)^\top
   &=\nabla_X\nabla_YZ-A_{B(Y,Z)}X,\\
 (\bar\nabla_X\bar\nabla_YZ)^\perp
   &=B(X,\nabla_YZ)+\nabla_X^\perp(B(Y,Z)).
\end{align*}
减去$X,Y$交换项及$\bar\nabla_{[X,Y]}Z$。切向结果为
\[
 (\bar R(X,Y)Z)^\top=R(X,Y)Z-A_{B(Y,Z)}X+A_{B(X,Z)}Y.
\]
移项并与$W$取内积即得Gauss式。法向结果减去两份$\nabla B$后，剩下
$B(\nabla_XY-\nabla_YX-[X,Y],Z)=0$，故得Codazzi式。

再对$\xi$使用Weingarten公式，两次求导的法向部分为
\[
 (\bar\nabla_X\bar\nabla_Y\xi)^\perp
 =-B(X,A_\xi Y)+\nabla_X^\perp\nabla_Y^\perp\xi.
\]
反对称化给
\[
 (\bar R(X,Y)\xi)^\perp
 =R^\perp(X,Y)\xi-B(X,A_\xi Y)+B(Y,A_\xi X).
\]
后两项与$\eta$配对为
$-g(A_\eta X,A_\xi Y)+g(A_\eta Y,A_\xi X)
=-g([A_\xi,A_\eta]X,Y)$，整理即得第三式。整个推导在任意标架成立，不要求预先选择主方向。
\end{proof}

\begin{corollary}[截面曲率与欧氏缩并]\label{cor:submanifold-curvature-traces}
若$u,v$正交归一，则
\[
 K_M(u,v)=\bar K(u,v)+\bar g(B(u,u),B(v,v))-|B(u,v)|^2.
\]
特别对欧氏浸入，令$|B|^2=\sum_{i,j}|B(e_i,e_j)|^2$，有
\[
 \Ric(X,X)=m\bar g(B(X,X),\mathbf H)-\sum_i|B(X,e_i)|^2,
 \qquad\mathrm{Scal}=m^2|\mathbf H|^2-|B|^2.
\]
因此欧氏极小子流形的$\Ric\leq0$、$\mathrm{Scal}=-|B|^2\leq0$；不由此断言其每个截面曲率均非正。
\end{corollary}
\begin{proof}
在Gauss式代入$(X,Y,Z,W)=(u,v,v,u)$即得首式。欧氏环境曲率为零，再将$(e_i,X,X,e_i)$代入并对$i$求和，第一项成为$\bar g(B(X,X),\sum_iB(e_i,e_i))$，第二项为所列平方和。再取$X=e_j$并求和，得到标量式。极小时去掉含$\mathbf H$项即可。维数大于2时，若干截面曲率的和非正，不足以逐项确定符号。
\end{proof}

\subsection*{超曲面、法平坦与环境曲率项}
在常截面曲率$c$的环境中，$\bar R(X,Y)Z=c(\bar g(Y,Z)X-\bar g(X,Z)Y)$。对超曲面，在一点取$S$的正交特征基，$Se_i=k_i e_i$，于是
\[
 K_M(e_i,e_j)=c+k_i k_j\quad(i\ne j),\qquad
 (\nabla_XS)Y=(\nabla_YS)X.
\]
后一式来自环境曲率的法向部分为零及$\nabla^\perp n=0$。一般环境中不能略去这一法向项。正球面外法向给$k_i=-1/a$，仍有$K=a^{-2}$；圆柱面一主曲率为零，故内蕴平坦。

在欧氏环境中，Ricci方程说明
\[
 R^\perp=0\quad\Longleftrightarrow\quad
 [A_\xi,A_\eta]=0\text{ 对所有法向量 }\xi,\eta.
\]
因为右式与所有$X,Y$配对为零等价于算子为零。有限组可交换自伴算子可同时正交对角化：先分解一个算子的特征空间，其余算子因交换保持各特征空间，再在这些空间内递归对角化。取法空间有限基即可应用。因此法平坦给逐点共同主方向；它不自动给一个处处光滑、甚至全局定义的主方向标架。超曲面的法联络局部单位标架中恒零，所以法曲率恒零；但全局法丛的平行绕行仍可能变号。

\Needspace{16\baselineskip}
\begin{proposition}[欧氏全脐超曲面的局部分类]\label{prop:umbilic-classification}
设$m\geq2$，欧氏超曲面在连通邻域选有单位法向$n$且$S=\lambda I$。则$\lambda$为常数；$\lambda=0$时像局部为仿射超平面的开集，$\lambda\ne0$时为半径$1/|\lambda|$球面的开集。未假设完备时不能将开集改为整个球面或平面。
\end{proposition}
\begin{proof}
Codazzi式给$(X\lambda)Y=(Y\lambda)X$。在任意点，给单位$X$选与之正交的单位$Y$，与$Y$取内积得$X\lambda=0$，故$\dd\lambda=0$。连通性给常数。

若$\lambda=0$，则$\dd n=0$，$n=n_0$恒定，且$\dd\langle f,n_0\rangle=0$，故像在某仿射超平面中。若$\lambda\ne0$，则
\[
 \dd(f+\lambda^{-1}n)=\dd f+\lambda^{-1}\dd n=0,
\]
所以$C=f+\lambda^{-1}n$恒定，$|f-C|=1/|\lambda|$。两种情形中，$f$到相应$m$维超曲面的微分均为同维同构，逆映射定理说明它在每点附近是到开集的微分同胚。维数1时$S=\lambda I$对每条曲线自动成立，却不迫使曲率恒定。
\end{proof}
